\documentclass[11pt,a4paper]{article}
\usepackage[T1]{fontenc}
\usepackage{lmodern,amsmath,amssymb,booktabs,microtype}
\usepackage[margin=24mm]{geometry}
\usepackage[colorlinks=true,urlcolor=blue,linkcolor=blue]{hyperref}
\newcommand{\EE}{\mathbb E}
\newcommand{\tr}{\operatorname{tr}}
\newcommand{\Rem}{\operatorname{Rem}}
\newcommand{\W}{\mathsf W_2}
\title{A positive residual in the quadratic memory sector\\Fixed-charge norm and the completed equal-time Gram}
\author{Alan Day \and Chengheng\thanks{AI research assistant (OpenAI).}}
\date{Local research note --- 6 November 2026}
\begin{document}\maketitle
\begin{abstract}
We evaluate the missing half-filled norm of the derivative singlet.
A Casimir identity for the exchanged contraction reduces the
adjacent-charge calculation to two known Racah kernels. An explicit
expansion cancels the terms of order $p^2$ and $p$; logarithmically
growing windows control the remainders and the complete input-rank
sum. Combined with the preceding cross, this proves a positive
output-deficit residual after projection onto the three existing
directions $H^2$, $A$, and $J_HF$.
\end{abstract}
\section{Setting and main result}
Use standard CAR, the full overlapping complex Gaussian $p$-body
ensemble, $\hbar=1$, and the normalized maximally mixed trace at
$m=N/2$. Inner products include the disorder average. Let
$h=H_m$, $k=H_{m-1}$ use the same couplings and set
\begin{equation}
 C=\EE\tr_m h^2,\quad B=\W h^2,\quad
 O=\W A,\quad A=\Rem(k^2)+(2p-m)h^2.
\end{equation}
$\W$ is ordinary Gaussian second chaos. The output-deficit operator
is $\mathcal S_N=\sum_{a\leq b,r}(a+b-r)\W Z^+_{ab;r}$,
with the rank-product convention of the preceding notes.
First take $p^2/N\to\lambda>0$, $C\to J_s^2/\lambda$;
write $q=e^{-4\lambda}$ and $\theta=2\lambda$. We prove
\begin{align}
 \frac{\|O\|_{2,m}^2}{C^2}&\longrightarrow
       4\lambda+4\lambda^2(1+q),\label{Aresult}\\
 \|O-\mathcal S_N\|_{2,m}^2&\longrightarrow
       \Delta_*:=\frac{2J_s^4(1+q)}\lambda>0.\label{difference}
\end{align}
Let $V_S=J_s^4[2(1-q)/\lambda+16q/(1+q)]$ be the
previous output-deficit variance after the $B$ subtraction.
Projection onto the span of $B,O,J_HF$ leaves
\begin{equation}
 V_{\rm new}=\frac{V_S\Delta_*}{V_S+\Delta_*}
 =\frac{J_s^4(1+q)(1-q^2+8\lambda q)}
 {\lambda(1+q+4\lambda q)}>0.\label{new}
\end{equation}
After this double-scaled limit, $\lambda\to0^+$ gives
$V_{\rm new}\to16J_s^4$. The result concerns this specified
equal-time operator span in the fixed half-filled state.

\newpage
\section{The exchanged Casimir insertion}
This step works at any charge $M$ away from the trivial endpoints.
For independent isotropic Hermitian rank-$a,b$ tensors $U,V$ with
$\EE U^2=\EE V^2=I$, define
\begin{equation}
 \mathcal C X=\frac12\sum_{ij}[E_{ij},[E_{ji},X]],\qquad
 \mathcal C U=x_aU,\quad x_a=a(N+1-a).
\end{equation}
Write $\Phi_a(X)=\EE UXU$, whose rank-$b$ eigenvalue is the
known neutral Racah kernel $\chi^{(M)}_{ab}$.
The ordered-product moment from the earlier covariance calculation
and the exchanged-product moment are
\begin{align}
 \EE\tr_M (UV)^\dagger\mathcal C(UV)
 &=x_a+x_b-\frac{x_ax_b}{C_f^{(M)}},\notag\\
 \EE\tr_M UV\mathcal C(UV)
 &=(x_a+x_b)\chi^{(M)}_{ab},\label{exchange}\\
 C_f^{(M)}&=\frac{M(N-M)(N+1)}N.\notag
\end{align}
For completeness, the second identity has a short direct proof.
The product rule gives
$\mathcal C(UV)=(x_a+x_b)UV+\sum_{ij}[E_{ij},U][E_{ji},V]$.
For a fixed pair $E=E_{ij},F=E_{ji}$, the extra trace, averaged
over $U$, is
\begin{equation}
 \tr_M\{[\Phi_a(VE)-\Phi_a(V)E][F,V]\}=0.
\end{equation}
Indeed $\Phi_a$ is self-adjoint for the trace pairing and
$\Phi_a([F,V])=\chi^{(M)}_{ab}[F,V]$ by equivariance.
The two terms cancel. The uninserted exchanged trace is
$\EE\tr UVUV=\chi^{(M)}_{ab}$. This is an application of
the established covariance decomposition \cite{kota}.

Return to $m=N/2$ and put $n=m+1$. The removal map satisfies
\begin{equation}
 \|\Rem X_r\|_{2,m}^2=
 \frac m n[mn-x_r]\|X_r\|_{2,m-1}^2.
\end{equation}
Let $C_a=\EE h_a^2$, $\omega_a=C_a/C$, and
\begin{equation}
 R_a=\frac{C_a^-}{C_a}=
 \frac{(m-p)^2 n}{m(mn-x_a)},\qquad
 C_-^f=\frac{(m-1)n(N+1)}N.
\end{equation}
Using \eqref{exchange}, the norm of the removed lower product,
per ordered input pair divided by $C_aC_b$, is
\begin{equation}
 R_aR_b\frac m n\left\{
 (mn-x_a-x_b)(1+\chi^-_{ab})+\frac{x_ax_b}{C_-^f}\right\}.
 \label{removed}
\end{equation}
Here $\chi^-_{ab}=\chi^{(m-1)}_{ab}$. Ordinary Gaussian
decoupling includes the $a=b$ coefficient in the ordered-pair sum;
no independent-disorder replacement of $A$ is made.

\newpage
\section{Finite formula and explicit cancellation}
Let $\chi=\chi^{(m)}_{ab}$ and $\xi$ be the adjacent-charge kernel
from the preceding note. Its lower parameters are $(1-m,-m,-N-1)$.
If $c_j$ is the $j$th summand of $\chi$, then
\begin{equation}
 \xi_j=\frac m{m-j}c_j,\qquad
 (\chi^-)_j=\left[1+\frac{j}{(m+1)(m-j)}\right]c_j.
\end{equation}
Thus $\chi^-=\chi+(\xi-\chi)/n$ exactly. Define
\begin{equation}
 g=\frac{(m-p)^2}{m},\quad h_0=2p-m,\quad \eta=\frac{p^2}{m},
 \quad P=\frac{g^2x_ax_b}{(mn-x_a)(mn-x_b)},\quad D=\xi-\chi.
\end{equation}
Combining \eqref{removed} with the known coherent cross and the
upper-product norm yields
\begin{align}
 \|O\|_{2,m}^2/C^2&=\sum_{a,b}\omega_a\omega_b f_{ab},\notag\\
 f_{ab}&=\eta^2(1+\chi)+P\left(\frac{mn}{C_-^f}-1-\chi\right)
       +\left(\frac{g^2-P}{n}+2h_0g\right)D.\label{finite}
\end{align}
Every $f_{ab}$ is a nonnegative quadratic-coefficient norm.
The individual terms in this formula have large cancellations.

Set $a=p-u$, $b=p-v$, $m=p^2/(2\lambda)$ for the expansion.
With $c_j^{(0)}=(-4\lambda)^j/j!$,
\begin{align}
 c_j&=c_j^{(0)}\left[1+\frac{L_1}{p}
                  +\frac{L_2+L_1^2/2}{p^2}+\cdots\right],\notag\\
 L_1&=-j(u+v+2\lambda)-j(j-1),\notag\\
 L_2&=-\frac j2(u^2+v^2)-\frac{j(j-1)}2(u+v)
       -\frac{j(j-1)(2j-1)}6\notag\\
 &\hspace{6mm}+\lambda j(u+v+1)+\frac{3\lambda j(j-1)}2-\lambda^2j.
\end{align}
Put $K=(g^2-P)/n+2h_0g$. The coefficient multiplying $c_j$
in \eqref{finite} is $\eta^2-P+Kj/(m-j)$, with expansion
\begin{align}
 p^2r_2+pr_1+r_0+O(p^{-1}),\qquad
 r_2&=-\frac{j+4\lambda}{\lambda},\notag\\
 r_1&=4(2j+2\lambda+u+v),\notag\\
 r_0&=-2j^2+(1-20\lambda)j+4\lambda(u+v+2)-4uv.
\end{align}
The unsigned part $\eta^2+P(mn/C_-^f-1)$ tends to
$4\lambda(1+\lambda)$. The three series coefficients sum to
\begin{equation}
 \sum_j c_j^{(0)}(r_2,\ r_1+r_2L_1,\
 r_0+r_1L_1+r_2[L_2+L_1^2/2])=(0,0,4\lambda^2q).
 \label{cancellation}
\end{equation}
These are polynomial identities: apply powers of $z\partial_z$
to $e^z$ and set $z=-4\lambda$. Both divergent orders cancel,
and the finite coefficient is independent of $u,v$.

\newpage
\section{Uniform control of the sums}
The expansion on page 3 must survive multiplication by $p^2$.
Here is a useful way to control the actual sums. Fix $\lambda$ in
a compact subset of $(0,\infty)$ and take
$K_p=\lceil\log p\rceil$.

For all $a,b\leq p$ and sufficiently large $p$, the exact term ratio
gives $|c_j|\leq A_\lambda^j/j!$. Moreover $m-j\geq m-p$.
The entire coefficient $\eta^2-P+Kj/(m-j)$ is bounded by
$B_\lambda p^2(1+j)$. Therefore the absolute contribution from
$j>K_p$, even before cancellation, tends to zero faster than any
inverse power of $p$.

On $u,v,j\leq K_p$, each logarithm used to expand $c_j$ has its
argument within $O(K_p/p)$ of one. Taylor's remainder bound,
summed over the $j$ factors, gives the displayed second-order
expansion with a relative remainder
\begin{equation}
 O\!\left(\frac{(1+u+v+j)^d}{p^3}\right)
\end{equation}
for a fixed finite degree $d$. The rational coefficients admit the
same kind of uniform expansion because their scaled denominators
are bounded away from zero. After multiplication by $p^2$ and
summation with the factorial envelope, the remainder is at most
$O((\log p)^{d'}/p)\to0$. Truncating the polynomial sums in
\eqref{cancellation} at $K_p$ has a superpolynomially small error.
Consequently
\begin{equation}
 \sup_{0\leq u,v\leq K_p}
 \left|f_{p-u,p-v}-[4\lambda+4\lambda^2(1+q)]\right|\to0.
\end{equation}

Finally, the previous input-rank ratio bound gives
$\omega_{p-u}\leq A^u/u!$ up to a constant depending on
$\lambda$. Outside this window it suffices to use the crude
positive bound $f_{ab}\leq B_\lambda p^4$. It follows either
from \eqref{removed} and Cauchy--Schwarz, or from
$\|\Rem\|\leq m$, bounded $R_a$, and the symmetric-product
norm bound. Hence the total contribution of $u>K_p$ or $v>K_p$
vanishes. This proves \eqref{Aresult} for the complete rank sum.
Using $\lambda_N=p^2/N$ in the same compact bounds covers sequences
with $\lambda_N\to\lambda$, not just exact rational sequences.

The argument also permits any fixed polynomial weight in $u,v$:
outside the window its extra factor is at most a power of $p$,
still dominated by the factorial tail. In particular the positive
$A$ norm is tight in its two input deficits. This result does not
control the output deficit $S=a+b-r$ of $A$; that separate tail is
still needed for the nonconstant memory filter.

\newpage
\section{The completed Gram and the surviving direction}
Order the operators as $(B,L_N,\mathcal S_N,O)$, where $L_N$
is the centered input-deficit quadratic operator from the preceding
notes. Put $z=1+q$, $b=\theta(1-q)$,
$s=\theta^2z+\theta(1-q)$, and $a=\theta^2z+2\theta$.
The old parity Gram, the preceding coherent cross, and
\eqref{Aresult} now give the full matrix
\begin{equation}
 C^{-2}G_N\longrightarrow
 \begin{pmatrix}
 z&0&b&b\\
 0&\theta z&0&0\\
 b&0&s&s\\
 b&0&s&a
 \end{pmatrix}.\label{gram}
\end{equation}
Define $v=s-b^2/z=\theta(1-q)+4\theta^2q/z$ and
$d=a-s=\theta z$. Its determinant is $\theta^2z^3v>0$.
In particular the fourth direction is not lost in double scaling.
The finite Gram of the three old directions has an invertible
limit, so its Schur complement converges to that of \eqref{gram}.

For an explicit limiting orthogonalization let
$c=b/z$, $Y=O-cB$, $\mathcal S_\perp=\mathcal S_N-cB$,
and $\rho=v/(v+d)$. Then
\begin{equation}
 T_N=\mathcal S_\perp-\rho Y
\end{equation}
has asymptotically zero inner product with $B,O,L_N$, and
$\|T_N\|^2\to C^2vd/(v+d)$, giving \eqref{new}.
The previously proved bridge
\begin{equation}
 J_HF-4\lambda q C L_N\longrightarrow0\quad\text{in }L^2
\end{equation}
can now also be paired with $O$, since \eqref{Aresult} supplies
a bounded fixed-sector norm. Thus the same conclusion holds for
the span $(B,O,J_HF)$.

Equation \eqref{difference} follows directly from
$\langle\mathcal S_N,O\rangle\to\|\mathcal S_N\|^2$.
It explains why the nonzero cross found in the preceding note does
not eliminate the output-deficit direction: $O-\mathcal S_N$
retains its own positive strength. The variance of $Y$ is now
\begin{equation}
 V_S+\Delta_*=4J_s^4\left[\lambda^{-1}+\frac{4q}{1+q}\right].
\end{equation}
It agrees with the earlier full-Fock expression, but here it was
derived in the fixed half-filled state rather than imported from
that expression.

This is linear independence in a specified equal-time span. A
complete projection onto all conserved operators and a real-time
spectral measure for $T_N$ have not been obtained. A new bulk
particle does not follow solely from this positive norm.

\newpage
\section{Evidence and next calculation}
The four saved computational records support distinct parts of the
derivation:
\begin{itemize}
 \item The general-charge Casimir identity is checked using all
 68,544 ordered Gaussian coordinate pairs over 15 rank-pair cases;
 the 45 exchanged, ordered, and uninserted identities are exact.
 \item The lower-charge contiguous relation has 384 finite exact
 checks. The raw and combined norm formulas agree in 1,200 cases.
 The latter also reproduce the three previous complete CAR norms
 $28$, $81$, and $2010$ at covariance $w=1$.
 \item Symbolic term-ratio checks verify $L_1,L_2$. The two
 divergent coefficients vanish identically in $\lambda,u,v$;
 the finite coefficient, Gram determinant, Schur complement,
 and small-$\lambda$ value $16$ are checked symbolically.
 \item Twelve all-rank calculations at $\lambda=0.1,0.25,0.5$
 and $p=16,32,64,128$ are evaluated at 80 and 120 digits.
 The finite $(B,\mathcal X_N/N,O)$ Gram is independently compared
 with all 27 entries of the three saved CAR Grams, and two
 further norms with exact rational pair sums.
\end{itemize}
The numerical difference norm and Schur complement use the
Casimir proxy $\mathcal X_N/N$, not the finite $\mathcal S_N$.
The preceding $L^2$ proxy bound and the now bounded $O$ norm
justify the same limiting targets. At $\lambda=0.5,p=128$ the
three quantities (divided by $C^2$) differ from their limits by
approximately $-0.0310954$, $-0.0126200$, and $-0.0261320$.
These finite-$p$ deviations are retained. Cancellation factors
and precision refinement are recorded; no interval-rounding
certificate is claimed.

The next missing estimate is the \emph{output}-rank tail of the
positive $A$ norm after the adjacent-charge cancellation. Together
with the input-tail result on page 4, it should permit the actual
memory coefficient $g_A(a,b;r)$ to be replaced by its fixed-deficit
limit in $L^2$. The central charge window and the low-chaos terms
of the fourth-order descendants must then be restored to assemble
the raw second memory coefficient. Finite-memory correlations and
the bulk space/mass dictionary remain subsequent tasks.

\begin{thebibliography}{9}
\bibitem{kota} V. K. B. Kota and M. Vyas,
\emph{Random Matrix Theory for Transition Strength Densities in
Finite Quantum Systems: Results from Embedded Unitary Ensembles},
\href{https://arxiv.org/abs/1504.00759}{arXiv:1504.00759},
Eqs. (8)--(14), (34)--(35), and Sec. 7.
The covariance and recoupling framework is used here; its dilute
asymptotics are not the half-filled limit of this note.
\end{thebibliography}
\end{document}
